\section{Vulkan and CUDA: existing assistance and a future hybrid}
\label{architecture-cuda-interop}

Nimbus can explore NVIDIA-specific computation without making NVIDIA hardware
a prerequisite for its portable solver. There are two different technical
routes: launching a limited PTX kernel from a Vulkan command buffer, which the
source already implements for retained acoustic flux, and sharing explicit
resources between Vulkan and the CUDA runtime, which would require additional
engineering. Neither route establishes an arbitrary CUDA solver backend or
multi-GPU domain decomposition.

Performance work should examine a specific operation and the workload around
it. CUDA may reduce the work required, provide a better implementation or make
a suitable library available while preserving Nimbus's numerical and ownership
contracts. Crossing an API boundary also adds dispatch and synchronization
costs. Complete-step measurements must establish whether the combined design
improves the actual simulation.

\subsection{The narrow path already present in Nimbus}
The strict-v1 selector \path{extensions.compute.cuda_assist} defaults to false
and requires the Vulkan backend. It does not add \texttt{cuda} to the supported
backend names. Its
\href{../vulkan/include/compute/vulkan_cuda_kernel.hpp}{kernel owner}
loads the retained-flux PTX module, obtains the function entry point and records
its dispatch inside a Vulkan command buffer using
\texttt{vkCmdCudaLaunchKernelNV}. The strict retained-flux path uses Vulkan
buffer device addresses for its density, wind and flux buffers. The rest of the
solver's scientific profile and native ownership requirements remain in force.

Initialization checks the selected NVIDIA device, Vulkan extension availability,
buffer-device-address support and the CUDA-launch feature. The current helper
requires compute capability 7.5 or newer. Missing header/API support, unavailable
entry points or module creation failure rejects that backend initialization.
The general runtime can disclose a CPU recovery when the profile permits
\path{extensions.compute.allow_fallback}; native-required execution must have
effective fallback disabled. Merely requesting assist is therefore not proof of
its execution. The active path logs the retained-flux PTX digest, and its manifest
identity includes that digest.

The readiness registry names this limited retained-flux capability separately.
Its source and bounded evidence are not a statement that every operator can run
through CUDA, that a long hardware campaign has passed, or that a broader hybrid
has completed scientific review. The
\href{https://docs.vulkan.org/refpages/latest/refpages/source/VK_NV_cuda_kernel_launch.html}{Khronos extension description}
explains that this mechanism dispatches PTX through Vulkan without creating the
usual Vulkan/CUDA context-interoperability path. It is a provisional,
NVIDIA-specific extension, so driver and header availability remain part of
admission and release maintenance.

\subsection{A different future route: shared external resources}
A more general hybrid could keep Vulkan as the owner of ordinary solver
dispatch and use a CUDA runtime stream for a selected operation or library.
The proposed first scope is one process, one physical NVIDIA GPU and explicitly
shared buffers. It is not implemented by the existing assist flag. A future
implementation would need a distinct, opt-in public capability and an exact
supported profile before users could select it as a normal execution path.

NVIDIA's
\href{https://docs.nvidia.com/cuda/cuda-programming-guide/04-special-topics/graphics-interop.html}{Vulkan interoperability guide}
describes importing exported memory and synchronization handles into CUDA on
the same physical device, identified by matching device UUIDs. An ordinal such
as device zero is insufficient because the two APIs can enumerate devices
differently. On a laptop, Studio's webview may render on another adapter from
the solver. This proposed sharing concerns solver allocations; it does not give
the separate Studio process direct access to them.

The aim would be to avoid a full device-to-host-to-device round trip around the
selected operation. The shared allocation remains GPU-accessible while the two
APIs exchange permission to use its current contents. ``Shared'' does not mean
that both may write simultaneously or that a completed host API call proves
the device work finished. A producer/consumer protocol is still necessary.

\begin{figure}[htbp]
\centering
\begin{tikzpicture}[
 box/.style={draw=NimbusBlue,rounded corners,fill=NimbusLight,align=center,
   text width=3.5cm,minimum height=1.25cm,font=\small},
 flow/.style={-{Latex[length=2mm]},thick,draw=NimbusBlue}]
\node[box] (vk) {Vulkan producer\\Produce, release, signal};
\node[box,right=0.8cm of vk] (cuda) {CUDA stream\\Wait, compute, signal};
\node[box,right=0.8cm of cuda] (next) {Vulkan consumer\\Wait, acquire, validate};
\draw[flow] (vk) -- (cuda);
\draw[flow] (cuda) -- (next);
\node[draw=NimbusBlue,fill=white,rounded corners,below=0.9cm of cuda,
  text width=11.6cm,minimum height=1.0cm,align=center,font=\small] (allocation)
  {One explicitly shareable device allocation, matched device UUID\\
   Stage $s$, generation $g$: declared ranges, precision and one current writer};
\draw[flow] (vk.south) -- (allocation.north west);
\draw[flow] (cuda.south) -- (allocation.north);
\draw[flow] (allocation.north east) -- (next.south);
\end{tikzpicture}
\caption{Proposed CUDA-runtime interoperability on one physical GPU. External
memory avoids a host round trip; external synchronization and ownership still
order access. This is distinct from Nimbus's existing Vulkan-command-buffer
PTX retained-flux path and from future multi-GPU domain decomposition.}
\end{figure}


\subsection{Admission must query the actual allocation contract}
A proposed interop session would first identify the same physical device in
both APIs and check the supported operating-system handle types. It would then
query the intended resource, not just the device's marketing identity.
Khronos's
\href{https://docs.vulkan.org/refpages/latest/refpages/source/VkPhysicalDeviceExternalBufferInfo.html}{external-buffer query}
includes the buffer's usage, creation flags and external handle type. The
\href{https://docs.vulkan.org/refpages/latest/refpages/source/VkExternalMemoryProperties.html}{returned memory properties}
describe export/import features and compatible handle types. A valid design
must honor dedicated-allocation requirements as well as ordinary alignment and
memory-type requirements.

For Nimbus, the imported object's record would need allocation size, field byte
range, element precision, strides, alignment, domain generation and owner. A
suballocation offset cannot be dropped simply because two APIs refer to the
same underlying memory. The CUDA import's dedicated-resource flag must agree
with how Vulkan actually allocated the resource. CUDA buffer mappings and
image mappings are different interfaces: a tiled image cannot be assumed to
have a linear row layout merely because its logical shape matches a field.
An initial buffer-only scope would reduce the number of layout assumptions
that need to be proved.

The
\href{https://docs.nvidia.com/cuda/cuda-runtime-api/cuda_runtime_api/group__CUDART__EXTRES__INTEROP.html}{CUDA external-resource API reference}
specifies mapping offsets and alignment, separate release of mapped buffers,
and handle ownership. Successful opaque-FD imports transfer descriptor ownership
to CUDA; ordinary Win32 NT handles retain application-managed handle lifetime.
Outstanding semaphore operations must finish before their objects are destroyed.
Nimbus would need separate unwind paths for failure before import, after import,
after mapping and after submitted work, so an error does not leak an allocation
or close a handle twice.

An external-memory handle also does not make host access automatically coherent.
If the selected policy maps the allocation on the CPU, the required noncoherent
flush/invalidate operations and barriers still apply. The runtime API documents
this requirement for CPU-mapped Vulkan memory. A host diagnostic or output
writer needs a published, completed endpoint, just as it does in the existing
Vulkan residency design.

\subsection{Ordering two APIs without serializing the whole device}
The proposed operation would have an explicit dependency chain: Vulkan produces
the input range and releases the required ownership; an exported semaphore
signals that endpoint; the CUDA stream waits, executes the selected operation
and signals completion; Vulkan waits and reacquires the resource before its
next consumer. The exact barriers, stage/access scopes and any image layouts
must follow the resource type and external-ownership rules. The
\href{https://docs.vulkan.org/spec/latest/chapters/synchronization.html}{Vulkan synchronization specification}
distinguishes execution order, memory visibility and queue-family ownership.
One of those requirements cannot stand in for all the others.
Where the sharing contract requires external ownership transfer, the barrier
uses the applicable external queue-family identity, such as
\texttt{VK\_QUEUE\_FAMILY\_EXTERNAL}; a CUDA device ordinal is not a Vulkan
queue-family index. Image layout and ownership must be tracked separately.

The
\href{https://docs.vulkan.org/refpages/latest/refpages/source/VkExternalSemaphoreProperties.html}{external-semaphore capability query}
must admit the selected semaphore and handle type. Binary semaphores need a
correct signal/wait reuse protocol; a timeline design additionally needs support
for its exported handle and a consistent monotonically advancing value scheme.
The CUDA guide describes both cases and their cross-API signal/wait ordering.
For Nimbus, a timeline value would be associated with a concrete field generation
and integrator stage rather than treated as an unexplained global counter.

CUDA's asynchronous external-semaphore signal is ordered in its specified
stream after earlier work in that stream. A library operation on another
stream is not thereby complete. The hybrid owner would set the library's stream
explicitly or insert the required cross-stream dependency before signaling.
Similarly, submitting a Vulkan command buffer does not establish that the
CUDA wait has completed. Full-device idle waits can provide a simple diagnostic
baseline, but using them around every operation could erase the benefit of
interoperability. Any later overlap optimization must retain the same dependency
graph and failure semantics.

\subsection{Which computations might justify the additional path?}
Candidate work should be chosen from measured solver costs. A substantial
well-defined numerical operation with stable input/output buffers is a better
initial experiment than repeatedly switching APIs around tiny pointwise kernels.
A batched operation can amortize launch and handoff overhead; a long sequence
of small transfers can instead increase latency and duplicate scratch storage.

NVIDIA provides
\href{https://docs.nvidia.com/cuda/cublas/index.html}{cuBLAS} for dense linear
algebra, \href{https://docs.nvidia.com/cuda/cusparse/index.html}{cuSPARSE} for
sparse operations and \href{https://docs.nvidia.com/cuda/cufft/index.html}{cuFFT}
for Fourier transforms. These are possible building blocks, not evidence that
Nimbus currently needs or calls them. Mapping a stencil, column solve or
diagnostic to a library routine is worthwhile only when its formulation,
boundary conditions, precision and data-layout costs fit the actual algorithm.
Introducing a spectral method, for example, changes more than the dispatch API
if the original operation was not mathematically spectral.

CUDA can also provide a useful development environment for a specialized
kernel. That does not automatically make its generated arithmetic faster on
the target GPU, nor remove the need to maintain the portable Vulkan, Metal and
CPU implementations. The maintenance cost includes another compiler/runtime
matrix, library versions, failure paths, scientific parity and hardware tests.
A future optional path should leave builds without the CUDA toolkit functional.

\subsection{Memory topology and multiple devices}
The benefit of avoiding host round trips can be especially relevant when a
discrete GPU has its own memory and a finite host interconnect. On a unified
system, copies and synchronization still have measurable costs, and the two API
contexts do not automatically share every allocation. Unified CPU/GPU memory,
CUDA managed memory and explicitly exported Vulkan memory are distinct
mechanisms. This NVIDIA-specific proposal does not establish CUDA support on
Apple GPUs or turn ARM64 into a proxy for unified-memory behavior.

Same-device resource sharing also does not establish peer access. Two installed
GPUs need their own device identities, import/transfer capabilities, synchronization
and ownership. An allocation imported for one device is not a promise that
another device can use it. The
\hyperref[architecture-multigpu]{multi-GPU proposal} adds physical-domain
partitioning, halo exchange, reductions and continuation ownership beyond this
single-device experiment. A CUDA-aware communication library would require
evidence for the actual imported allocation and transport; its name alone does
not validate arbitrary Vulkan buffers.

\subsection{Scientific and failure-atomicity requirements}
A future hybrid operation would consume the same stage, carrier, physical domain,
boundary state and private moments as the existing owner. It would produce the
same declared outputs and accounting. Duplicating only the main visible field
while leaving a private reservoir or diagnostic on an old endpoint would not
be an implementation of the original operation.

Precision must be explicit across the boundary. NVIDIA's
\href{https://docs.nvidia.com/cuda/floating-point/index.html}{floating-point
white paper} explains why fused operations, intermediate rounding and reduction
order can change results even with IEEE arithmetic. Nimbus would retain its
declared storage and intermediate types, FMA contract, denormal behavior and
reduction requirements. Library math modes or reduced-precision accelerators
could not be enabled merely to obtain a faster timing. Any intentional numerical
change would need its own scientific justification and tests.

Tests would include direct references, limiting and positivity, conservation or
closure, invalid inputs, threading behavior and CPU/Vulkan/Metal parity where
applicable to the declared profile. They would also exercise an error after
Vulkan production but before CUDA completion, and an error before Vulkan
reacquisition. A rejected transaction must leave the previous committed state,
private fields, reports and generation ownership consistent. Restart must
capture the full continuation endpoint after both APIs finish, not a snapshot
taken between them.

The public configuration, catalog, readiness record and runtime dispatch report
would need to identify the composite path and its exact supported scope.
Dispatch evidence should name the selected operation, device UUID, kernel or
library identity, precision, shared/staged transfer path and synchronization
outcome. Unsupported imports or capabilities should reject an explicitly
required hybrid profile. An intentionally permitted fallback must be disclosed;
it cannot count as a successful native hybrid experiment.

\subsection{An experiment needs a complete cost model}
For a proposed operation, an illustrative break-even condition is
\begin{equation}
 T_{\mathrm{Vulkan\ operation}} >
 T_{\mathrm{CUDA\ operation}}+T_{\mathrm{handoff}}+
 T_{\mathrm{layout}}+T_{\mathrm{extra\ synchronization}}.
\end{equation}
The right-hand terms include real preparation and ownership work. Persistent
mapping can amortize setup over many steps, but its retained memory and teardown
still count. A steady-state measurement that excludes compilation or import
must report those startup costs separately.

The bounded performance gate should compare the same source/configuration,
scientific endpoint, grid and output cadence on the same hardware. Record
compiler flags, driver and library versions, clock/power conditions, CPU worker
settings, memory topology, warmup and repeated-trial variation. Measure the
selected operation, complete solver step, host synchronization, output/restart
cost and peak memory. A faster kernel that moves the bottleneck to publication
or increases memory pressure can lose end to end.

Only after ordinary public selection, admission, engineering and bounded tests
are complete would this new capability be Ready for its required long campaigns.
Campaign validation and independent scientific review remain separate steps
toward Fully Qualified. The present design analysis supplies a path to test,
not a speedup promise or a declaration that external-resource interop is shipped.
